\documentclass[nm]{../../elegantbook}
% ========== Хамар хэсэг болон суурь мэдээлэл ==========
\title{ElegantBook тестийн баримт}
\subtitle{Загварын үндсэн шинж чанарууд}
\author{Тестийн инженер}
\institute{Загварын тестийн баг}
\date{\today}
\version{v1.0}
\extrainfo{Энэ баримт зөвхөн загварын функцийн туршилтанд зориулагдсан, ямар ч судлаах үнэ цэнэгүй}
\bioinfo{Тестийн багц}{TEST-FULL-20260506}
\logo{../../figure/logo-blue.png}
\cover{../../figure/cover.jpg}

% ========== Гарчиг болон суурь тохируулга ==========
\setcounter{tocdepth}{3} % Гарчиг 3 түвшин хүртэл харуулна

% ========== Тестийн санал баримт ==========
\addbibresource[location=local]{../en/reference.bib}

\begin{document}

\maketitle
\frontmatter
\tableofcontents

% ========== Удиртгал (Section-тэй хувааж) ==========
\pagenumbering{Roman}
\setcounter{page}{0}
\chapter*{Тестийн удиртгал}
\markboth{Тестийн удиртгал}{}
\section*{Баримтын тайлбар}
\markright{Баримтын тайлбар}
Энэ баримт бол ElegantBook бүх функцийг турших загвар бөгөөд бүх контент бол хуурамч тестийн өгөгдөл бөгөөд реальный утгагүй.
\newpage
\section*{Тестийн хүрээ}
\markright{Тестийн хүрээ}
Энэ тест нь дараахь зүйлсийг агуулна: хамар хэсэг, гарчиг, үсэг, олон түвшний жагсаалт, төвөг томъёо, бүх математик орчин, хөрвүүлэлт, зураг/хүснэгт, хавсаргах тайлбар, бүлгийн товчлол, дасгал, санал баримт, хавсаргах хэсэг.

\mainmatter

% ========== 1 Бүлэг: Базар байршуулалт болон олон түвшний жагсаалт ==========
\chapter{Базар байршуулалт болон жагсаалтын туршилт}
\begin{introduction}[Энэ бүлгийн тестийн зүйлс]
\item Монгол үсэг болон текстын формат
\item 3 түвшний тэгш жагсаалтын гүнзгийрүүлэлт
\item 3 түвшний дараалсан жагсаалтын гүнзгийрүүлэлт
\item Хэвийн параграф болон тодотгосон текст
\end{introduction}

\section{Үсэг болон текстын туршилт}
Хэвийн основ текст. \textbf{Тодотгосон}、\textit{Налуу}、\underline{Доогуур шугам}、\textcolor{red}{Улаан текстын туршилт}。

\section{3 түвшний тэгш жагсаалтын туршилт}
\begin{itemize}
\item Нэгдүгээр түвшний тест A
\item Нэгдүгээр түвшний тест B
    \begin{itemize}
    \item Хоёрдугаар түвшний тест B1
    \item Хоёрдугаар түвшний тест B2
        \begin{itemize}
        \item Гуравдугаар түвшний тест B2-1
        \item Гуравдугаар түвшний тест B2-2
            \begin{itemize}
            \item Дөрөвдүгээр түвшний тест B2-1a
            \item Дөрөвдүгээр түвшний тест B2-2b
            \end{itemize}
        \end{itemize}
    \end{itemize}
\end{itemize}

\section{3 түвшний дараалсан жагсаалтын туршилт}
\begin{enumerate}
\item Нэгдүгээр түвшний дараалсан тест 1
\item Нэгдүгээр түвшний дараалсан тест 2
    \begin{enumerate}
    \item Хоёрдугаар түвшний дараалсан тест 2-1
    \item Хоёрдугаар түвшний дараалсан тест 2-2
        \begin{enumerate}
        \item Гуравдугаар түвшний дараалсан тест 2-2-1
        \item Гуравдугаар түвшний дараалсан тест 2-2-2
            \begin{enumerate}
            \item Дөрөвдүгээр түвшний дараалсан тест 2-2-2-1
            \item Дөрөвдүгээр түвшний дараалсан тест 2-2-2-2
            \end{enumerate}
        \end{enumerate}
    \end{enumerate}
\end{enumerate}

% ========== 2 Бүлэг: Төвөг математик томъёоны туршилт ==========
\chapter{Төвөг математик томъёоны туршилт}
\begin{introduction}
\item Бодит мөр төвөг томъёо
\item Олон мөртэй томъёо (align)
\item Матриц/детерминант
\item Давхар интеграл/цуврал/хязгаар/фракц
\item Томъёоны шошго болон хөрвүүлэлт
\end{introduction}

\section{Бодит мөр төвөг томъёо}
Мөр дэх төвөг томъёо：$\sum_{i=1}^n \frac{x_i^2}{\sqrt{y_i+1}} \geq \bar{x}^2$、$\lim_{n\to\infty} \left(1+\frac{1}{n}\right)^n = e$。

\section{Олон мөртэй томъёо болон матриц}
\begin{align}
a^2 + b^2 &= c^2 \label{eq:gougu} \\
\int_{0}^{1}\int_{0}^{1} f(x,y) dxdy &= 1 \label{eq:double} \\
\begin{vmatrix}
a & b \\
c & d
\end{vmatrix} &= ad-bc \label{eq:det}
\end{align}

\section{Том оператортой томъёо}
\begin{equation}\label{eq:series}
\sum_{k=0}^{\infty} \frac{x^k}{k!} = e^x,\quad \oint_{C} Pdx+Qdy = \iint_{D}\left(\frac{\partial Q}{\partial x}-\frac{\partial P}{\partial y}\right)dxdy
\end{equation}

Томъёоны хөрвүүлэлт：Гуго теорем\eqref{eq:gougu}、Давхар интеграл\eqref{eq:double}、Детерминант\eqref{eq:det}、Цуврал\eqref{eq:series}。

% ========== 3 Бүлэг: Бүх математик орчны туршилт ==========
\chapter{Бүх математик орчны туршилт}
\begin{introduction}
\item Тодорхойлолт/томьёо/лемма/үр дүн/аксиом/постулат
\item Таамаглал/жишээ/бодлого/дасгал
\item Тэмдэглэл/тайлбар/дүгнэлт/урьдчилсан нөхцөл/шинж чанар
\item Хариулт/баталгааны орчин
\end{introduction}

\section{Орчны туршилт}

% 1. Тодорхойлолт
\ifdefined\simpleTest
\begin{definition}[Тестийн тодорхойлолт]\label{def:test}
\else
\begin{definition}[Тестийн тодорхойлолт]{test}
\fi
Энэ бол тодорхойлолтын орчны туршилтын контент бөгөөд ямар ч математик утгагүй, зөвхөн загварыг шалгахад зориулагдсан.
\end{definition}

% 2. Томьёо
\ifdefined\simpleTest
\begin{theorem}[Тестийн томьёо]\label{thm:test}
\else
\begin{theorem}[Тестийн томьёо]{test}
\fi
Хэрэв $x>0$ бол $x+1>1$ бөгөөд \eqref{eq:gougu} томъёотой холбоотой.
\end{theorem}

% 3. Лемма
\ifdefined\simpleTest
\begin{lemma}[Тестийн лемма]\label{lem:test}
\else
\begin{lemma}[Тестийн лемма]{test}
\fi
Бүх бодит тоо $x$-д $x^2\geq0$ байдаг.
\end{lemma}

% 4. Үр дүн
\ifdefined\simpleTest
\begin{corollary}[Тестийн үр дүн]\label{cor:test}
\else
\begin{corollary}[Тестийн үр дүн]{test}
\fi
Лемма\ref{lem:test}-ээс дараахь үр дүн гарна：$(x-1)^2\geq0$。
\end{corollary}

% 5. Аксиом
\ifdefined\simpleTest
\begin{axiom}[Тестийн аксиом]\label{axi:test}
\else
\begin{axiom}[Тестийн аксиом]{test}
\fi
Тестийн аксиом бөгөөд ямар ч утгагүй.
\end{axiom}

% 6. Постулат
\ifdefined\simpleTest
\begin{postulate}[Тестийн постулат]\label{pos:test}
\else
\begin{postulate}[Тестийн постулат]{test}
\fi
Тестийн постулат бөгөөд зөвхөн орчныг шалгахад зориулагдсан.
\end{postulate}

% 7. Таамаглал
\ifdefined\simpleTest
\begin{proposition}[Тестийн таамаглал]\label{pro:test}
\else
\begin{proposition}[Тестийн таамаглал]{test}
\fi
Тестийн таамаглалын контент бөгөөд тодорхойлолт\ref{def:test}-тэй холбоотой.
\end{proposition}

% 8. Жишээ
\begin{example}\label{exa:test}
Жишээ орчны туршилт, автоматаар дугаарлагдана.
\end{example}

% 9. Бодлого
\begin{problem}\label{probl:test}
Бодлогын орчны туршилт, шийдэх хэрэгтэй бодлого гаргахад зориулагдсан.
\end{problem}

% 10. Дасгал
\begin{exercise}\label{exer:test}
Тооцоол：$1+2+3+\dots+10$。
\end{exercise}

% 11. Тэмдэглэл
\begin{note}
Тэмдэглэл орчин：дугааргүй, чухал сануулгад зориулагдсан.
\end{note}

% 12. Тайлбар
\begin{remark}
Тайлбар орчин：нэмэлт тайлбарлахад зориулагдсан.
\end{remark}

% 13. Дүгнэлт
\begin{conclusion}
Дүгнэлт орчин：тогтмол дүгнэлтийн контент.
\end{conclusion}

% 14. Урьдчилсан нөхцөл
\begin{assumption}
Урьдчилсан нөхцөл орчин：шалтгааны тохируулгалт.
\end{assumption}

% 15. Шинж чанар
\begin{property}
Шинж чанар орчин：обьектийн шинж чанар.
\end{property}

% 16. Хариулт
\begin{solution}
Дасгал\ref{exer:test} хариулт：Нийт бол $55$ (тестийн хариулт)。
\end{solution}

% 17. Баталгаа
\begin{proof}
Томьёо\ref{thm:test} баталгаа：$x>0$ тул хоёуланд нь 1 нэмэхээр баталгаажна.
\end{proof}

% ========== 4 Бүлэг: Хөрвүүлэлт·Зураг/Хүснэгт·Хавсаргах тайлбар ==========
\chapter{Хөрвүүлэлт болон нэмэлт функц}
\begin{introduction}
\item Бүх төрлийн хөрвүүлэлт
\item Энгийн зураг/хүснэгтийн туршилт
\item Хавсаргах тайлбарын функц
\item Бүлгийн шилжилтийн шалгалт
\end{introduction}

\section{Бүх хөрвүүлэлтийн цуглуулга}
Тодорхойлолт\ref{def:test}、Томьёо\ref{thm:test}、Лемма\ref{lem:test}、Үр дүн\ref{cor:test}、Аксиом\ref{axi:test}、Постулат\ref{pos:test}、Таамаглал\ref{pro:test}、Жишээ\ref{exa:test}、Бодлого\ref{probl:test}、Дасгал\ref{exer:test}。

\section{Зураг/хүснэгтийн туршилт}
\begin{figure}[h]
    \centering
    \includegraphics[width=0.5\textwidth]{../../image/scatter.jpg}
    \caption{Тестийн зураг}\label{fig:test}
\end{figure}
\begin{table}[h]
    \centering
    \begin{tabular}{c*{6}{c}}
    \toprule
        $n$ & 1 & 2 & 3 & 4 & 5 & 6 \\
    \midrule
        $a_n$ & 1 & 1 & 2 & 3 & 5 & 8 \\ 
    \bottomrule
    \end{tabular}
    \caption{Тестийн хүснэгт}\label{tab:test}
\end{table}
Зураг/хүснэгтийн хөрвүүлэлт：\figref{fig:test} бол тестийн зураг，\tabref{tab:test} бол тестийн хүснэгт。

% ========== Бүлгийн хойш дасгал ==========
\begin{problemset}[Энэ бүлгийн тестийн дасгалууд]
\item Тодорхойлолт\ref{def:test}ийн харагдах загварыг шалгах
\item \eqref{eq:double} томъёог гаргах
\item Зураг\ref{fig:test} зөв харуулж байгаа эсэхийг шалгах
\end{problemset}

% ========== Санал баримт ==========
\nocite{*}
\printbibliography

% ========== Хавсаргах хэсэг ==========
\appendix
\chapter{Тестийн хавсаргах хэсэг}
Хавсаргах хэсгийн контент нь хавсаргах бүлэг, толгой хэсэг, хуудасны дугаар, байршуулалтын загварыг шалгахад зориулагдсан бөгөөд ямар ч мэдээлэлгүй.
\section{Хавсаргах туршилт 1}
Хавсаргах хэсгийн контент нь хавсаргах бүлэг, толгой хэсэг, хуудасны дугаар, байршуулалтын загварыг шалгахад зориулагдсан бөгөөд ямар ч мэдээлэлгүй.
\newpage
\section{Хавсаргах туршилт 2}
Хавсаргах хэсгийн контент нь хавсаргах бүлэг, толгой хэсэг, хуудасны дугаар, байршуулалтын загварыг шалгахад зориулагдсан бөгөөд ямар ч мэдээлэлгүй.
\end{document}